\documentclass[10pt,a4paper]{article}
\usepackage[T1]{fontenc}
\usepackage[utf8]{inputenc}
\usepackage{lmodern}
\usepackage[margin=21mm]{geometry}
\usepackage{amsmath,amssymb,booktabs,tabularx,array,longtable,xcolor,hyperref,fancyhdr,graphicx,enumitem,float}
\hypersetup{colorlinks=true,urlcolor=blue!55!black,linkcolor=blue!55!black}
\setlength{\parindent}{0pt}\setlength{\parskip}{.48em}
\setlength{\emergencystretch}{3em}
\setlist[enumerate]{leftmargin=1.5em,itemsep=.2em}
\pagestyle{fancy}\fancyhf{}
\fancyhead[L]{\small PCM / Tune-1 ternary search}
\fancyhead[R]{\small 9 November 2026}
\fancyfoot[C]{\thepage}\setlength{\headheight}{14pt}
\begin{document}
{\LARGE\bfseries PCM Tune-1: bounded ternary hyperparameter search}\\[5pt]
{\large 20\% sampling, NPL-default noise, 256-by-256 images}\\[4pt]
{\small Preparation report, version 1.0}\\[4pt]
\input{status.tex}
\par\medskip\hrule

\section{Executive summary}
We prepare an experiment with two independent PnP algorithms:
PnP-ADMM-DRUNet and GS-PnP-GS-DRUNet.
\emph{Test set 450} remains the full original test set
(90 images in each of five categories).
\emph{Test set 45} contains nine randomly drawn held-out test
images per category, sampled with a fixed random seed.
The 45 image IDs and their original measurement hashes are frozen.

We select \textbf{five validation calibration images per category}.
Each image produces one independent tune-1 search.
There are therefore five category experiments containing
\textbf{25 independent searches for each PnP algorithm}, or
50 searches altogether. Each resulting configuration is applied
unchanged to the same nine held-out test images of its category.
Parameters are selected only from noisy acquired measurement
cross-validation; test PSNR, SSIM and image ground truth are
never used for tuning.

\section{Tuning algorithm}
For an acquisition with Hadamard model $y=Ax+\varepsilon$,
let $\Lambda_k$ be one of four pre-frozen held-out sets of acquired
fine-scale Hadamard rows, and let $\Theta_k$ be its complement
within the acquired rows. The first 4,096 coarse patterns always
stay in the fit set. The original 9,011 other acquired patterns
are divided among four held-out folds.
For candidate $\theta$, calculate
\[
L(\theta)=\frac14\sum_{k=1}^4
\frac{\|A_{\Lambda_k}\widehat x_\theta(A_{\Theta_k},y_{\Theta_k})
-y_{\Lambda_k}\|_2^2}{|\Lambda_k|}.
\]
Each image has its own loss and independently selected
configuration. Every candidate is evaluated using the \emph{same}
four measurements splits for that calibration image.

\subsection*{Log-coordinate ternary search}
The previous 15- and 18-candidate grids were Cartesian products of
specific parameter values. The new strategy instead specifies
continuous positive ranges and shrinks them numerically.
For a single parameter, let $u=\log_{10}p$ and start with
$[a,b]=[\log_{10}p_{\min},\log_{10}p_{\max}]$.
At each step, evaluate two interior points,
\[
u_1=a+(b-a)/3,\qquad u_2=b-(b-a)/3.
\]
Under the assumption that the one-dimensional loss slice
is \emph{unimodal}, retain $[a,u_2]$ if
$L(u_1)\le L(u_2)$, otherwise retain $[u_1,b]$.

With several parameters, hold the other coordinates fixed,
perform three ternary contractions for one coordinate,
then move to the next. Use two passes, reversing the
coordinate order in the second pass. Keep the best
\emph{measured, numerically valid} configuration.
One-dimensional unimodality does \textbf{not} imply a globally
unimodal two- or three-dimensional objective.
Discontinuous solver stopping and backtracking can violate
this assumption in practice.

\begin{center}
\begin{tabular}{@{}lll@{}}
\toprule
Algorithm & Parameter & Frozen bounds\\
\midrule
PnP-ADMM & $\rho$ & 0.03 to 4.0\\
PnP-ADMM & Denoiser $\sigma$ & 0.015 to 0.18\\
GS-PnP & $\lambda$ & 0.005 to 0.40\\
GS-PnP & Denoiser $\sigma$ & 0.015 to 0.15\\
GS-PnP & Initial $\tau_0$ & 3.0 to 75\\
\bottomrule
\end{tabular}
\end{center}

Use the original pre-trained denoisers, normalised PCM
fidelity and solver stopping rules.
Every candidate is scored on all four SSDU folds.
Numerical failures invalidate a candidate; they must be
recorded rather than assigned a fictitious finite score.

\section{Step-by-step plan}
\begin{enumerate}
\item \textbf{Freeze sample selection.}
Select five validation and nine test images per category,
independently of reconstruction metrics. Save RNG seeds,
image IDs and measurement/mask SHA256 fingerprints.
\textbf{Completed.}
\item \textbf{Verify search code.}
Check interval shrinking on synthetic unimodal objectives,
reproducibility, output hashes and test-access guards.
\textbf{Completed in the CPU preparation tests.}
\item \textbf{Perform 50 tune-1 searches.}
Run 25 ADMM and 25 GS-PnP independent searches on the
preselected noisy validation data. Save all probed
parameters, four fold losses, solver diagnostics,
and failed configurations.
\item \textbf{Freeze all parameters before test evaluation.}
Test-set access is blocked until every one of the 50
independent selection records has been frozen.
\item \textbf{Apply parameters to test45.}
For each calibration, reconstruct each of the nine
test images in the same category, using the full
20\% acquired measurement set. This produces
225 full reconstructions per PnP family.
\item \textbf{Compare against established methods.}
Reuse archived 450-image baseline metrics after checking
case and acquisition hashes. Regenerate any missing
historical image arrays under the original fixed settings,
then verify their image metrics.
\item \textbf{Aggregate correctly.}
For each tune-1 calibration, calculate the mean and
sample SD over nine test images. To report one
test45 number for tune-1, average five calibration
results for each \emph{unique} test image first,
then average across the 45 unique test images.
The repeated calibrations are not 225 independent
test-image samples.
\item \textbf{Prepare deliverables.}
Save individual reconstructions and scores, paired
comparisons for PSNR, SSIM, HaarPSI and MSE, per-category
tables and plots, and 45 visual comparison pages
showing all historical and new tune-1 results.
\end{enumerate}

\section{Experiment running}
\textbf{Frozen image IDs (per category):}
\input{selection_table.tex}

Both PnP solvers use the same 20\% Hadamard sampling
rate and fixed NPL-default noise already archived
for validation and test. No noise or masks will be
regenerated for parameter selection.
The planned GPUs are tuni RTX 4070 (ADMM) and
RTX 4090 (GS-PnP). The prepared runner isolates
each GPU and saves scores after each calibration
candidate, allowing safe resumption.

Two parameter passes with three ternary contractions
per coordinate require at most 25 distinct candidate
evaluations per ADMM calibration and 37 per GS calibration.
At four folds per candidate, upper budgets are
$25\times25\times4=2{,}500$ ADMM solver calls and
$25\times37\times4=3{,}700$ GS solver calls.
A solver call can take up to 100 iterations and
GS-PnP may require additional network calls during
backtracking. These are \emph{planned budgets}, not
completed computations.

\section{Results and deliverables}
\input{results_section.tex}

The preparation archive includes the frozen test45
and validation sample lists, parameter bounds,
tested pure ternary-search code, validation-only
GPU tuning runner, test45 reconstruction runner,
historical baseline replay scripts, metric
aggregator and qualitative gallery generator.
Full numerical results will only be inserted
after all searches and reconstructions finish.

\section{Conclusion}
\input{conclusion_dynamic.tex}

\paragraph{Primary sources.}
Yaman et al. (2020), \emph{Self-Supervised Learning of Physics-Guided
Reconstruction Neural Networks without Fully-Sampled Reference Data},
section on measured-data holdout:
\url{https://doi.org/10.1002/mrm.28378}.

Berkeley Statistics 243, \emph{Univariate Function Optimisation},
section on unimodal search and interval bracketing:
\url{https://stat243.berkeley.edu/fall-2024/units/unit11-optim.html}.

Hurault, Leclaire and Papadakis (2022),
\emph{Gradient Step Denoiser for Convergent Plug-and-Play},
GS-PnP update and backtracking:
\url{https://openreview.net/pdf?id=fPhKeld3Okz}.
\end{document}